% ==============================================================================
% Resumo em língua estrangeira (obrigatório, NBR 14724:2024)
% ==============================================================================
\begin{resumo}[Abstract]
\begin{instrucao}
\orientacao{Abstract}
Versão do resumo em inglês, com o mesmo conteúdo e as mesmas regras. Não use
tradução automática sem revisão cuidadosa. Termos técnicos devem seguir o uso
da área (por exemplo, \emph{user story} para ``história de usuário'', e não uma
tradução literal). Se usar ferramenta de IA para traduzir ou revisar, declare
o uso no apêndice sobre inteligência artificial.
\end{instrucao}

\begin{otherlanguage*}{english}
Academic tutoring is one of the main forms of student support in the initial
courses of Computing programs, but scheduling sessions usually relies on
informal messages and spreadsheets, which leads to schedule conflicts and
makes it difficult for coordinators to follow up. This work presents Monitora,
a web system for scheduling tutoring sessions that allows tutors to publish
their available hours, students to book those hours and professors to track
attendance and the topics covered. The system was developed with an iterative
process based on Scrum, in six two-week cycles. Requirements were gathered in
interviews with tutors and professors and specified as user stories. The
implementation uses the Django framework, the PostgreSQL database and Docker
containers. Verification combined automated tests, with 87\% code coverage,
and an evaluation with 18 users, who gave the system an average score of 84.6
on the System Usability Scale, close to the range rated as excellent. The results indicate that
the system meets the gathered requirements and can replace the informal
process currently used in the program's courses.

\keywords{web system; scheduling; academic tutoring; software engineering}
\end{otherlanguage*}
\end{resumo}
